\section{USAMO 1978}
        \begin{enumerate}
        	\item (\textit{USAMO 1978}) Given that $a,b,c,d,e$ are real numbers such that


 $a+b+c+d+e=8$,


 $a^2+b^2+c^2+d^2+e^2=16$.


 Determine the maximum value of $e$.


 



\textbf{Solution 1}

By Cauchy Schwarz, we can see that $(1+1+1+1)(a^2+b^2+c^2+d^2)\geq (a+b+c+d)^2$thus $4(16-e^2)\geq (8-e)^2$ Finally, $e(5e-16) \geq 0$ which means $\frac{16}{5} \geq e \geq 0$so the maximum value of $e$ is $\frac{16}{5}$.

 \textbf{from:} \href{https://artofproblemsolving.comhttp://image.ohozaa.com/view2/vUGiXdRQdAPyw036}{Image from Gon Mathcenter.net}

 



\textbf{Solution 2}

Seeing as we have an inequality with constraints, we can use Lagrange multipliers to solve this problem.We get the following equations:

 $(1)\hspace*{0.5cm} a+b+c+d+e=8\\ (2)\hspace*{0.5cm} a^{2}+b^{2}+c^{2}+d^{2}+e^{2}=16\\ (3)\hspace*{0.5cm} 0=\lambda+2a\mu\\ (4)\hspace*{0.5cm} 0=\lambda+2b\mu\\ (5)\hspace*{0.5cm} 0=\lambda+2c\mu\\ (6)\hspace*{0.5cm} 0=\lambda+2d\mu\\ (7)\hspace*{0.5cm} 1=\lambda+2e\mu$

 If $\mu=0$, then $\lambda=0$ according to $(6)$ and $\lambda=1$ according to $(7)$, so $\mu \neq 0$. Setting the right sides of $(3)$ and $(4)$ equal yields $\lambda+2a \mu= \lambda+2b \mu \implies 2a\mu=2b \mu \implies a=b$. Similar steps yield that $a=b=c=d$. Thus, $(1)$ becomes $4d+e=8$ and $(2)$ becomes $4d^{2}+e^{2}=16$. Solving the system yields $e=0,\frac{16}{5}$, so the maximum possible value of $e$ is $\frac{16}{5}$.

 



\textbf{Solution 3}

A re-writing of Solution 1 to avoid the use of Cauchy Schwarz. We have
 \[(a+b+c+d)^2=(8-e)^2,\] and 
 \[a^2+b^2+c^2+d^2=16-e^2.\]The second equation times 4, then minus the first equation, 
 \[(a-b)^2+(a-c)^2+(a-d)^2+(b-c)^2+(b-d)^2+(c-d)^2=4(16-e^2)-(8-e)^2.\]The rest follows.

 J.Z.

 



\textbf{Solution 4}

By the \href{https://artofproblemsolving.com/wiki/index.php?title=Principle\_of\_Insufficient\_Reasons}{Principle of Insufficient Reasons}, since $a,b,c,d$ are indistinguishable variables, the maximum of $e$ is acheived when $a=b=c=d$, so we have 
 \[4a+\max e=8\] 
 \[4a^2+(\max e)^2=16\] 
 \[\implies \max e=\boxed{\frac{16}{5}}\]. $\square$ ~\href{https://artofproblemsolving.com/wiki/index.php?title=Ddk001}{Ddk001}

 
                \begin{itemize}
                    \item Note: For some reason I think this solution is missing something.
                \end{itemize}


                


	\item (\textit{USAMO 1978}) $ABCD$ and $A'B'C'D'$ are square maps of the same region, drawn to different scales and superimposed as shown in the figure. Prove that there is only one point $O$ on the small map that lies directly over point $O'$ of the large map such that $O$ and $O'$ each represent the same place of the country. Also, give a Euclidean construction (straight edge and compass) for $O$.


 
                    \begin{figure}[H]
                        \centering
                        \includegraphics[width=0.5\textwidth]{Problem/USAMO_I_1978_2_Problem_Figure_1.png}
                    \end{figure}
                    


 



\textbf{Solution 1}

The point is obviously unique, because the two maps have different scales (but if P and Q where two fixed points the distance between them would be the same on both maps).

 Let the small map square be A'B'C'D' and the large be ABCD, where X and X' are corresponding points. We deal first with the special case where A'B' is parallel to AB. In this case let AA' and BB' meet at O. Then triangles OAB and OA'B' are similar, so O must represent the same point. So assume A'B' is not parallel to AB.

 Let the lines A'B' and AB meet at W, the lines B'C' and BC meet at X, the lines C'D' and CD meet at Y, and the lines D'A' and DA meet at Z. We claim that the segments WY and XZ meet at a point O inside the smaller square. W cannot lie between A' and B' (or one of the vertices A', B' of the smaller square would lie outside the larger square). If it lies on the opposite side of A' to B', then Y must lie on the opposite side of C' to D'. Thus the segment WY must cut the side A'D' at some point Z' and the side B'C' at some point X'. The same conclusion holds if W lies on the opposite side of B' to A', because then Y must lie on the opposite side of D' to C'. Similarly, the segment XZ must cut the side A'B' at some point W' and the side C'D' at some point Y'. But now the segments X'Z' and W'Y' join pairs of points on opposite sides of the small square and so they must meet at some point O inside the small square.

 Now the triangles WOW' and YOY' are similar (WW' and YY' are parallel). Hence OW/OY = OW'/OY'. So if we set up coordinate systems with AB as the x-axis and AD as the y-axis (for the large square) and A'B' as the x'-axis and A'D' as the y'-axis (for the small square) so that corresponding points have the same coordinates, then the y coordinate of O equals the y' coordinate of O. Similarly, XOX' and ZOZ' are similar, so OX/OZ = OX'/OZ', so the x-coordinate of O equals its x'-coordinate. In other words, O represents the same point on both maps.

 



\textbf{Solution 2}

The point is obviously unique because of the different scale and angle, which can also be seen as the point can and must be the unique center of spiral similarity taking ABCD to A'B'C'D', as under the spiral similarity the center does not move. (Recall that a spiral similarity is a dilation combined with a rotation, and that there is a unique spiral similarity taking any segment to another segment, provided the four points do not make a parallelogram.)To construct this center of spiral similarity, reference EGMO Lemma 10.1. Draw line DD' and AA', and let them intersect at point X. Draw circles through points X,A,D and X,A',D' (using the perpendicular bisector constructions to find the center of the circle using two chords). The intersection point of the two circles other than X is the center of spiral similarity, and hence is O, which represents the same point on both maps.

 


	\item (\textit{USAMO 1978}) An integer $n$ will be called \textit{good} if we can write


 $n=a_1+a_2+\cdots+a_k$,


 where $a_1,a_2, \ldots, a_k$ are positive integers (not necessarily distinct) satisfying


 $\frac{1}{a_1}+\frac{1}{a_2}+\cdots+\frac{1}{a_k}=1$.


 Given the information that the integers 33 through 73 are good, prove that every integer $\ge 33$ is good.


 



\textbf{Solution}

\textit{Lemma:} If $g$ is a good integer, then so is $2g + 2$ and $2g + 9$.

 Proof: Let $g = a_1 + a_2 + \cdots + a_k$ such that 
 \[\frac{1}{a_1}+\frac{1}{a_2}+\cdots+\frac{1}{a_k}=1,\] where the $a_i$ are positive integers. Then, notice that 
 \[1 = \frac{1}{2} + \frac{1}{2} = \frac{1}{2} + \frac{1}{2a_1}+\frac{1}{2a_2}+\cdots+\frac{1}{2a_k},\] so $g_1 = 2 + 2a_1 + 2a_2 + ... + 2a_k = 2 + 2g$ is also a good integer. Furthermore, 
 \[1 = \frac{1}{3} + \frac{1}{6} + \frac{1}{2} = \frac{1}{3} + \frac{1}{6} + \frac{1}{2a_1}+\frac{1}{2a_2}+\cdots+\frac{1}{2a_k},\] so $g_2 = 3 + 6 + 2a_1 + 2a_2 + ... + a_k = 9 + 2g$ is also good.

 This completes the lemma. Now, we use induction to show that the integers from 33 to $n$ are all good, where $n \ge 33$.

 Base case: Our given takes care of $33 \le n \le 73$.For the inductive step, we assume that the integers $k$ between $33$ and $73$, inclusive, are all good. Now, we do some casework on the parity of $k+1$.

 If $k+1$ is odd, then write $k+1 = 2g + 9$ for some integer $g$.

 It follows that because $k \ge 74$, we must have

 
 \[g = \frac{k}{2} - 4 \ge 37 - 4 = 33.\]

 Thus, $33 \le g < \frac{k}{2}$, and so $g$ is good by the inductive hypothesis. Applying the lemma gives $2g + 9 = k+1$ as a good integer.

 If $k+1$ is even, write $k+1 = 2g+2$ for some integer $g$. Therefore, 
 \[g = \frac{k-1}{2} > \frac{73-1}{2} = 36 > 33,\] so that $33 < g < \frac{k}{2}$, meaning that $g$ is a good integer by the inductive hypothesis. From the lemma, $k+1$ must be good as well.

 This completes the inductive step and thus the induction. Therefore, all integers $n \ge 33$ are good integers.

 


	\item (\textit{USAMO 1978}) (a) Prove that if the six dihedral (i.e. angles between pairs of faces) of a given tetrahedron are congruent, then the tetrahedron is regular.


 (b) Is a tetrahedron necessarily regular if five dihedral angles are congruent?


 



\textbf{Solution}

(a) Let $ABCD$ be the said tetrahedron, and let the inscribed sphere of $ABCD$ touch the faces at $W, X, Y, Z$. Then, $OW, OX, OY, OZ$ are normals to the respective faces. We know that the angle between any two normals is equal, so we have $|OW|=|OX|=|OY|=|OZ|$ at equal angles. Now, since 
 \[WX=2\cdot OW\sin\frac{\angle{WOX}}{2},\] and similar for the other sides, we have that $WXYZ$ is a regular tetrahedron. Now, the faces of $ABCD$ are the tangent planes at $W$, $X$, $Y$, and $Z$. Then, consider a $120^{\circ}$ rotation about $OW$. The rotation sends $X\mapsto Y$, $Y\mapsto Z$, and $Z\mapsto X$. Thus we have $AB=AC$, $AC=AD$, $AD=AB$, and $BC=CD=DB$. Performing rotations about the other axes yields that $ABCD$ has equal edges, so it is regular.

 (b) Consider the normals $OW, OX, OY, OZ$. We can perform a transformation in which we slightly shift $X$, $Y$, and $Z$ closer such that $\triangle XYZ$ is still equilateral, and such that all angles between every pair of normals is less that $180^{\circ}$. Then, shift $W$ such that $WX=WY=XY$. Then five of the distances are equal but the sixth is not.

 


	\item (\textit{USAMO 1978}) Nine mathematicians meet at an international conference and discover that among any three of them, at least two speak a common language. If each of the mathematicians speak at most three languages, prove that there are at least three of the mathematicians who can speak the same language.


 



\textbf{Solution}

Suppose no three mathematicians speak the same language. Then person A can share some language with at most 3 other delegates, because if he shared some language with 4 delegates there would be 3 with the same language. So there are 5 delegates who do not share a language with A. Let one of them be B. Using the same logic, one the remaining 4, let it be C, does not share a language with B. Now A, B, C do not have any common languages. Contradiction.

 


\end{enumerate}